\chapter{Colectomy}
\section*{Overview}
Colectomy removes part or all of the colon and reconnects the bowel or creates a temporary or permanent ostomy.
\section*{Why This Test or Procedure Is Ordered}
It may be required for colorectal cancer, diverticular disease, inflammatory bowel disease, obstruction, perforation, ischemia, or severe bleeding.
\section*{How This Test or Procedure Is Performed}
The operation may be open, laparoscopic, or robotic. The diseased segment is removed, lymph nodes may be sampled, and the bowel ends are joined when safe.
\section*{Risks and Recovery}
Leak, infection, bleeding, ileus, bowel obstruction, blood clots, injury to nearby organs, and need for an ostomy are possible. Recovery includes gradual diet advancement and monitoring of bowel function.
\section*{References}
\begin{enumerate}\item MedlinePlus. Colon Resection. U.S. National Library of Medicine; 2025.\item Current Medical Diagnosis and Treatment. McGraw-Hill; 2025.\end{enumerate}
\clearpage
